% ECONOMICS_PROBLEM: {"id":"C73-15","title":"Payoff characterization of convergence for a fixed mirror-descent rule","jel_code":"C73","source_id":"OP-0148"}
% FIRST_POSTED: 2026-10-09
% ATTEMPT_STATUS: unresolved
% ENTRY_KIND: research_attempt
\documentclass[11pt]{article}
\usepackage{amsmath,amssymb,amsthm,geometry,hyperref}
\geometry{margin=1in}
\newtheorem{theorem}{Theorem}
\newtheorem{proposition}[theorem]{Proposition}
\newtheorem{lemma}[theorem]{Lemma}
\newtheorem{corollary}[theorem]{Corollary}
\theoremstyle{definition}
\newtheorem{example}[theorem]{Example}
\newtheorem{remark}[theorem]{Remark}
\newcommand{\R}{\mathbb R}
\newcommand{\N}{\mathbb N}
\DeclareMathOperator{\NE}{NE}
\DeclareMathOperator{\cav}{cav}
\DeclareMathOperator{\supp}{supp}
\title{Payoff characterization of convergence for a fixed mirror-descent rule\\ GPT 6 Astra Ultra}
\author{Research attempt}
\date{9 October 2026}
\begin{document}
\maketitle

\section*{Target and chosen specialization}
The full target fixes a mirror-descent rule, including its regularizers and step sizes, and asks which payoff arrays make every interior trajectory approach the Nash set. This attempt treats the negative-entropy regularizer
\[
 h_i(x_i)=\sum_a x_{i,a}\log x_{i,a},\qquad
 x_{i,a}^{t+1}=\frac{x_{i,a}^t\exp(\eta_t v_{i,a}(x^t))}
 {\sum_b x_{i,b}^t\exp(\eta_t v_{i,b}(x^t))}.
\]
The formula follows by differentiating the strictly concave constrained objective; positivity and normalization give the unique optimizer. All initial probabilities are strictly positive. The complete classification below for summable step sizes is an elementary consequence of bounded total log-odds movement. Its priority is not claimed. The general agenda and the need to specify the dynamics come from \cite{potential}.

\section*{An exact classification for finite total learning time}
Call a game strategically trivial if, for every player $i$, every two of its actions $a,b$, and every pure opponent profile $s_{-i}$,
\[
 u_i(a,s_{-i})=u_i(b,s_{-i}).
\]
\begin{theorem}[Summable entropy steps]
Fix finite nonempty action sets and positive step sizes with $H=\sum_{t=1}^{\infty}\eta_t<\infty$. For the simultaneous entropy update, the following are equivalent:
\begin{enumerate}
\item Every interior initial profile yields distance to the Nash set tending to zero.
\item The game is strategically trivial.
\end{enumerate}
Every trajectory in every finite game converges to an interior profile under these step sizes, whether or not that profile is Nash.
\end{theorem}
\begin{proof}
Choose $M>0$ bounding all absolute pure payoffs. For any two actions,
\[
 \log\frac{x_{i,a}^{t+1}}{x_{i,b}^{t+1}}
 =\log\frac{x_{i,a}^{t}}{x_{i,b}^{t}}
 +\eta_t\bigl(v_{i,a}(x^t)-v_{i,b}(x^t)\bigr).
\]
The increments are absolutely summable, being bounded by $2M\eta_t$. All log odds therefore converge to finite real numbers. Fixing one reference action for each player and normalizing their limiting odds proves convergence to a profile with all coordinates strictly positive.

If the game is strategically trivial, all action payoffs of each player coincide at every mixed opponent profile. Every update is the identity, and every profile is Nash.

Conversely suppose that, at a pure opponent profile $s_{-i}$, some actions $a,b$ satisfy $u_i(a,s_{-i})-u_i(b,s_{-i})=g>0$. Choose $\theta>0$ so small that $4M\theta<g/2$. Select interior initial beliefs for every opponent $j\ne i$ such that
\[
 \sum_{c\ne s_j}\frac{x_{j,c}^1}{x_{j,s_j}^1}
 <\frac{\theta e^{-2MH}}{\max\{1,n-1\}}.
\]
The log-odds bound shows that the same sum of odds at any time is less than $\theta/\max\{1,n-1\}$. The probability that some opponent deviates from $s_{-i}$ under the product distribution is consequently less than $\theta$. The expected payoff difference between $a$ and $b$ differs from $g$ by at most $4M\theta$, since the difference takes values in $[-2M,2M]$. It is thus greater than $g/2$ at every time and at the limiting profile. Player $i$'s limiting probability of $b$ is positive. A profile giving positive mass to a strictly suboptimal action is not Nash. For $n=1$ the same argument uses no opponents and the payoff difference is exactly $g$.

The Nash set is closed: it is specified by finitely many weak continuous best-response inequalities. Since the trajectory converges to a point outside this set, its distance to the set cannot tend to zero. This disproves the universal-initialization property whenever the game is not strategically trivial.
\end{proof}

\section*{Infinite total time does not suffice}
\begin{proposition}[Matching pennies obstruction for every positive schedule]
Consider the two-player zero-sum game with row payoff matrix
$\left(\begin{smallmatrix}1&-1\\-1&1\end{smallmatrix}\right)$.
For the entropy update with any positive step-size sequence, no interior trajectory starting away from the unique Nash equilibrium approaches the Nash set.
\end{proposition}
\begin{proof}
Let $p,q$ be the probabilities of the first actions and put
$z=\log(p/(1-p))$, $w=\log(q/(1-q))$. The update is
\[
 z'=z+2\eta\tanh(w/2),\qquad
 w'=w-2\eta\tanh(z/2).
\]
Define
\[F(z,w)=\log\cosh(z/2)+\log\cosh(w/2).\]
This is strictly convex, with gradient
\[\nabla F(z,w)=\tfrac12(\tanh(z/2),\tanh(w/2)).\]
The inner product of this gradient with $(z'-z,w'-w)$ is zero. Strict convexity therefore gives $F(z',w')>F(z,w)$ unless $(z,w)=(0,0)$. Since $F(0,0)=0$ and $F$ is positive elsewhere, a trajectory starting elsewhere cannot tend to $(0,0)$. Direct best-response comparison gives the unique Nash equilibrium $p=q=1/2$. Convergence in distance to that singleton would force $z,w\to0$, a contradiction.
\end{proof}

\begin{proposition}[An exact potential game with a discrete two-cycle]
Let both players receive one when their binary actions differ and zero when they coincide. For a constant entropy step size $\eta>4$, some interior initial profile has an exact period-two orbit disjoint from the Nash set.
\end{proposition}
\begin{proof}
Identical payoffs make the common payoff an exact potential: each unilateral payoff difference equals the potential difference. In log odds the update is
\[
 z'=z-\eta\tanh(w/2),\qquad
 w'=w-\eta\tanh(z/2).
\]
There is $c>0$ with $\eta\tanh(c/2)=2c$. Indeed, the left side minus $2c$ is zero at zero, has derivative $\eta/2-2>0$ there, and is negative for $c>\eta/2$. The intermediate value theorem gives a positive root. Starting at $(z,w)=(c,c)$ gives $(-c,-c)$, and the next step returns to $(c,c)$. At either point both players put positive mass on both actions, but each strictly prefers the less likely opposing action. Neither profile is Nash. Since this is a finite orbit disjoint from the closed Nash set, its distance from that set has a positive minimum.
\end{proof}

\section*{What this settles and what it leaves open}
The summable-step theorem is a complete payoff characterization for one specified family of rules. The two adversarial examples show why neither divergence of total step size nor exact potential structure can alone classify arbitrary simultaneous discrete updates. They do not contradict a continuous-time convergence theorem or a theorem requiring sufficiently small steps. General nonsummable entropy schedules and other regularizers remain outside the classification proved here.

\begin{thebibliography}{9}
\bibitem{potential} M. Bichler, D. Legacci, P. Mertikopoulos, M. Oberlechner, and B. Pradelski, \emph{Characterizing the Convergence of Game Dynamics via Potentialness}, arXiv:2503.16285v1 (2025), Sections 3--4 and 6. \url{https://arxiv.org/html/2503.16285v1}. The primary preprint and its stated classification agenda were inspected on 9 October 2026. The discrete calculations here are independent proofs, not claims that its convergence results apply to every step size.
\end{thebibliography}

\section{Completion Estimate}
25\%

\end{document}
